% ═══════════════════════════════════════════════════════════════
%  chapters/ch04-protocol.tex
%  Протокол обмена: кадры, CRC16, команды
% ═══════════════════════════════════════════════════════════════

\chapter{Протокол обмена}
\label{ch:protocol}

\epigraph{<<Протокол~--- это язык, на котором говорят устройства.
Если язык плохой~--- они не поймут друг друга>>.}{}

\section{Введение}
\label{sec:protocol-intro}

В \texttt{ch03} мы увидели \emph{три мира}: осциллограф,
генератор и ПК. Они \emph{общаются} друг с другом. Но
\emph{как} именно? По \textbf{протоколу}.

\begin{ideabox}
\textbf{Протокол~--- это \emph{контракт}.} Он говорит:
<<если ты пришлёшь \emph{такие} байты, я отвечу \emph{такими}>>.
Без контракта~--- \emph{хаос}. С контрактом~--- \emph{порядок}.
\end{ideabox}

\subsection{Зачем нужен протокол}

Представьте, что ПК отправляет генератору команду
<<поставь частоту 1 кГц>>. Что именно \emph{уйдёт} в порт?
Один байт? Строка <<1000>>? А если генератор \emph{не
успел} прочитать~--- как ПК узнает, что команда \emph{потерялась}?

\textbf{Протокол} отвечает на все эти вопросы:
\begin{itemize}
    \item \textbf{Как} упаковать команду (формат).
    \item \textbf{Как} проверить целостность (CRC).
    \item \textbf{Как} понять, что это \emph{новый} кадр, а не
          \emph{продолжение} старого (sync).
    \item \textbf{Как} сопоставить запрос и ответ (seq).
\end{itemize}

\subsection{Что вы узнаете}

\begin{itemize}
    \item Из каких \emph{полей} состоит кадр.
    \item Что такое \texttt{sync}, \texttt{seq}, \texttt{crc16}.
    \item Как \emph{считать} CRC16/IBM.
    \item Какие \emph{команды} понимают генератор и осциллограф.
    \item Как \emph{устроен} кадр данных \texttt{OSC\_DATA}.
\end{itemize}

\section{Формат кадра}
\label{sec:protocol-frame}

\textbf{Кадр}~--- это \emph{единица} обмена. Каждая команда,
каждый ответ, каждый блок данных~--- это \emph{один} кадр.

\begin{figure}[htbp]
\centering
\adjustbox{max width=\textwidth}{%
\begin{tikzpicture}[
    field/.style={
        draw, thick,
        minimum height=1cm,
        align=center,
        font=\footnotesize\ttfamily,
        fill=white
    },
    header/.style={field, fill=clrUIData!20},
    payload/.style={field, fill=clrLogicData!20},
    crc/.style={field, fill=clrWarn!20},
    label/.style={
        font=\footnotesize,
        above,
        inner sep=2pt
    },
    size/.style={
        font=\scriptsize,
        below,
        inner sep=2pt
    }
]

% Поля кадра
\node[header, minimum width=1.5cm] (sync) at (0, 0) {sync};
\node[header, minimum width=1cm, right=0cm of sync] (ver) {ver};
\node[header, minimum width=1.5cm, right=0cm of ver] (seq) {seq};
\node[header, minimum width=1.2cm, right=0cm of seq] (cmd) {cmd};
\node[header, minimum width=1.5cm, right=0cm of cmd] (len) {len};
\node[payload, minimum width=3cm, right=0cm of len] (payload) {payload};
\node[crc, minimum width=1.5cm, right=0cm of payload] (crc) {crc16};

% Подписи сверху
\node[label] at (sync.north) {0xAA55};
\node[label] at (ver.north) {0x01};
\node[label] at (seq.north) {seq++};
\node[label] at (cmd.north) {cmd};
\node[label] at (len.north) {len};
\node[label] at (payload.north) {данные};
\node[label] at (crc.north) {CRC-16/IBM};

% Размеры снизу
\node[size] at (sync.south) {2 байта};
\node[size] at (ver.south) {1};
\node[size] at (seq.south) {2};
\node[size] at (cmd.south) {1};
\node[size] at (len.south) {2};
\node[size] at (payload.south) {len байт};
\node[size] at (crc.south) {2};

\end{tikzpicture}
}% end adjustbox
\caption{Структура кадра протокола. Сверху~--- значения,
снизу~--- размеры полей.}
\label{fig:protocol-frame}
\end{figure}

\subsection{Поле \texttt{sync}: маркер начала}

\begin{center}
\texttt{0x55 0xAA} (\texttt{0xAA55} в little-endian)
\end{center}

\textbf{Sync}~--- это <<якорь>>. Приёмник \emph{сканирует}
поток байтов, пока не найдёт \texttt{0x55}, за которым
следует \texttt{0xAA}. Всё, что было \emph{до}~--- \emph{мусор}.

\begin{keybox}
\textbf{Зачем нужен sync?} Порт~--- это \emph{поток байтов}.
Если приёмник \emph{пропустил} начало кадра (например,
из-за \emph{переполнения} буфера), он должен \emph{найти}
следующий. Sync~--- \emph{уникальная} последовательность,
которой \emph{нет} в данных.
\end{keybox}

\subsection{Поле \texttt{ver}: версия протокола}

\begin{center}
\texttt{0x01}
\end{center}

\textbf{Версия}~--- это \emph{запас} на будущее. Если
мы \emph{изменим} протокол (добавим поля)~--- версия
\emph{позволит} старым устройствам \emph{игнорировать}
новые кадры.

\subsection{Поле \texttt{seq}: счётчик}

\begin{center}
\texttt{0x0000}--\texttt{0xFFFF} (2 байта)
\end{center}

\textbf{Seq} (sequence number)~--- это \emph{счётчик}
запросов. Каждый новый кадр \emph{увеличивает} его на 1.
Ответ содержит \emph{тот же} \texttt{seq}.

\begin{ideabox}
\textbf{Зачем нужен seq?} Если ПК отправил \emph{две}
команды подряд, а генератор \emph{ответил} только
на \emph{одну}~--- по \texttt{seq} можно понять,
\emph{какая} потерялась. Также \texttt{seq} \emph{защищает}
от \emph{дубликатов} (если кадр пришёл дважды~---
второй \emph{игнорируется}).
\end{ideabox}

\subsection{Поле \texttt{cmd}: код команды}

\begin{center}
\texttt{0x10}--\texttt{0x1F} (генератор), \texttt{0x20}--\texttt{0x2F} (осциллограф)
\end{center}

\textbf{Cmd}~--- это <<что делать>>. Диапазоны:
\begin{itemize}
    \item \texttt{0x10}--\texttt{0x1F}~--- команды генератору.
    \item \texttt{0x20}--\texttt{0x2F}~--- команды осциллографу.
    \item \texttt{0x40}~--- кадр данных \texttt{OSC\_DATA}.
    \item \texttt{0x80}~--- бит ответа (например, \texttt{0x90}~--- ответ на \texttt{0x10}).
\end{itemize}

\subsection{Поле \texttt{len}: длина данных}

\begin{center}
\texttt{0x0000}--\texttt{0xFFFF} (2 байта)
\end{center}

\textbf{Len}~--- \emph{сколько} байт в \texttt{payload}.
Нужна, чтобы \emph{знать}, когда кадр \emph{закончится}.

\subsection{Поле \texttt{payload}: данные}

\begin{center}
\texttt{len} байт
\end{center}

\textbf{Payload}~--- \emph{полезная нагрузка}. Для команды
\texttt{set\_freq}~--- это \SI{4}{\byte} (частота).
Для \texttt{OSC\_DATA}~--- это \emph{семплы} сигнала.

\subsection{Поле \texttt{crc16}: контрольная сумма}

\begin{center}
\texttt{0x0000}--\texttt{0xFFFF} (2 байта)
\end{center}

\textbf{CRC16}~--- это \emph{контрольная сумма}. Считается
от поля \texttt{ver} до конца \texttt{payload}. Если
при вычислении \emph{совпала} с принятой~--- кадр \emph{целый}.
Если \emph{нет}~--- кадр \emph{битый}.

\begin{warnbox}
\textbf{CRC не защищает от всего.} Она \emph{ловит}
\emph{случайные} ошибки (помехи, сбой USB). Но
\emph{не} ловит \emph{умышленные} подмены. Для
наших задач~--- \emph{достаточно}.
\end{warnbox}

\section{CRC16/IBM: как считается}
\label{sec:protocol-crc16}

\textbf{CRC16/IBM} (также известный как \textbf{CRC-16-ANSI})
~--- это \emph{стандартный} алгоритм. Полином~---
\texttt{0xA001}, начальное значение~--- \texttt{0xFFFF}.

\begin{listing}[htbp]
\begin{minted}[frame=none, bgcolor=gray!5]{c}
uint16_t crc16_ibm(const uint8_t *data, size_t len)
{
    uint16_t crc = 0xFFFF;

    for (size_t i = 0; i < len; i++) {
        crc ^= data[i];
        for (uint8_t bit = 0; bit < 8; bit++) {
            if (crc & 1) {
                crc = (crc >> 1) ^ 0xA001;
            } else {
                crc >>= 1;
            }
        }
    }
    return crc;
}
\end{minted}
\caption{\texttt{crc16\_ibm()}~--- подсчёт CRC16}
\label{lst:protocol-crc16}
\end{listing}

\subsection{Что такое XOR (исключающее ИЛИ)}

Прежде чем разбирать алгоритм~--- поймём \emph{одну}
операцию. \textbf{XOR} (в C~--- \texttt{\^{}})~--- это
\emph{побитовое} <<исключающее ИЛИ>>. Оно возвращает
\texttt{1}, \emph{только} если биты \emph{разные}:

\begin{center}
\begin{tabular}{cc|c}
\textbf{A} & \textbf{B} & \textbf{A \texttt{\^{}} B} \\
\hline
0 & 0 & 0 \\
0 & 1 & 1 \\
1 & 0 & 1 \\
1 & 1 & 0 \\
\end{tabular}
\end{center}

Например, \texttt{0xAA \^{} 0x0F}:
\[
\texttt{10101010} \oplus \texttt{00001111} = \texttt{10100101}
\]

\subsection{Как работает алгоритм}

\begin{enumerate}
    \item Начинаем с \texttt{crc = 0xFFFF}.
    \item Для \emph{каждого} байта:
    \begin{itemize}
        \item \emph{Выполняем} \texttt{crc \^{}= data[i]}~---
              \emph{XOR} младшего байта \texttt{crc} с байтом данных.
        \item Для \emph{каждого} из \emph{восьми} бит:
        \begin{itemize}
            \item Если младший бит \texttt{crc} равен \texttt{1}~---
                  \emph{сдвигаем} вправо на 1 и \emph{XOR} с \texttt{0xA001}.
            \item Если \texttt{0}~--- просто \emph{сдвигаем} вправо.
        \end{itemize}
    \end{itemize}
    \item Возвращаем \texttt{crc}.
\end{enumerate}

\begin{keybox}
\textbf{Почему именно \texttt{0xA001}?} Это \emph{обратный}
полином \texttt{0x8005}. Число \texttt{0xA001}~---
это \emph{стандарт} для CRC16/IBM. \emph{Менять} его
нельзя~--- иначе другая сторона \emph{не} поймёт
контрольную сумму.
\end{keybox}

\begin{ideabox}
\textbf{Зачем так сложно?} CRC~--- это \emph{математически
обоснованный} способ \emph{поймать} ошибки. Простая
сумма байтов \emph{не} ловит \emph{перестановки}.
CRC ловит \emph{большинство} случайных искажений~---
и \emph{очень} быстрый (всего \emph{несколько} операций
на байт).
\end{ideabox}



\begin{keybox}
\textbf{Проверка CRC~--- это \emph{фильтр}.} Если CRC \emph{не}
совпала~--- кадр \emph{выбрасываем} и \emph{продолжаем}
сканировать поток. Это \emph{защищает} от мусора.
\end{keybox}

\section{Команды генератора}
\label{sec:protocol-gen-cmds}

Генератор понимает \emph{шесть} команд (и \emph{одну}~--- статус).

\begin{table}[htbp]
\centering
\caption{Команды генератора (0x10–0x1F)}
\label{tab:protocol-gen-cmds}
\small
\begin{tabularx}{\textwidth}{@{}c l l L@{}}
\toprule
\textbf{Код} & \textbf{Имя} & \textbf{Payload} & \textbf{Назначение} \\
\midrule
\texttt{0x10} & set\_wave   & \texttt{u8 type} & Форма сигнала \\
\texttt{0x11} & set\_freq   & \texttt{u32 mHz} & Частота (миллигерцы) \\
\texttt{0x12} & set\_ampl   & \texttt{u16 mVpp} & Амплитуда \\
\texttt{0x13} & set\_offset & \texttt{i16 mV} & Смещение \\
\texttt{0x14} & set\_duty   & \texttt{u16 permille} & Скважность \\
\texttt{0x15} & upload\_wave & \texttt{u16 n, n×u16} & Пользовательская форма \\
\texttt{0x1F} & get\_gen\_status & --- & Запрос статуса \\
\bottomrule
\end{tabularx}
\end{table}

\subsection{Формы сигнала (0x10)}

\texttt{u8 type}:
\begin{itemize}
    \item \texttt{0}~--- синус.
    \item \texttt{1}~--- двухполупериодный.
    \item \texttt{2}~--- однополупериодный.
    \item \texttt{3}~--- пила.
    \item \texttt{4}~--- треугольник.
    \item \texttt{5}~--- меандр.
    \item \texttt{6}~--- пользовательская.
\end{itemize}

\subsection{Частота (0x11)}

\texttt{u32 mHz}~--- частота в \emph{миллигерцах}
(1 Гц = 1000 мГц). \textbf{Зачем так?} Чтобы \emph{избежать}
дробных чисел. \SI{1.5}{\kilo\hertz} = \num{1500000} мГц.

\subsection{Ответ (0x1F)}

Ответ на \texttt{get\_gen\_status}~--- это кадр с \texttt{cmd=0x9F}:

\begin{minted}[frame=none, bgcolor=gray!5]{text}
u8 wave; u32 freq_mHz; u16 ampl_mVpp;
i16 offset_mV; u16 duty_permille;
\end{minted}

\section{Команды осциллографа}
\label{sec:protocol-osc-cmds}

Осциллограф понимает \emph{пять} команд (и \emph{одну}~--- статус).

\begin{table}[htbp]
\centering
\caption{Команды осциллографа (0x20–0x2F)}
\label{tab:protocol-osc-cmds}
\small
\begin{tabularx}{\textwidth}{@{}c l l L@{}}
\toprule
\textbf{Код} & \textbf{Имя} & \textbf{Payload} & \textbf{Назначение} \\
\midrule
\texttt{0x20} & set\_fs      & \texttt{u32 Hz} & Частота дискретизации \\
\texttt{0x21} & set\_gain    & \texttt{u8 step} & Усиление \\
\texttt{0x22} & set\_trigger & \texttt{u8 mode, i16 level\_mV, u8 edge} & Триггер \\
\texttt{0x23} & capture\_once & \texttt{u16 pre\_pct, u16 samples} & Одиночный захват \\
\texttt{0x24} & stream\_on   & \texttt{u8 on} & Запуск/остановка потока \\
\texttt{0x2F} & get\_osc\_status & --- & Запрос статуса \\
\bottomrule
\end{tabularx}
\end{table}

\subsection{Триггер (0x22)}

\texttt{set\_trigger}~--- \emph{самая сложная} команда:
\begin{itemize}
    \item \texttt{mode}~--- \texttt{0} выкл, \texttt{1} нормальный, \texttt{2} авто.
    \item \texttt{level\_mV}~--- уровень, в милливольтах.
    \item \texttt{edge}~--- \texttt{0} по фронту, \texttt{1} по спаду.
\end{itemize}

\subsection{Поток (0x24)}

\texttt{stream\_on}~--- \emph{ключевая} команда. Включает
\emph{непрерывную} отправку кадров \texttt{OSC\_DATA}.





%%%%%=======================================




\section{Кадр данных \texttt{OSC\_DATA}}
\label{sec:protocol-osc-data}

\texttt{OSC\_DATA}~--- это \emph{особый} кадр. В нём
\texttt{cmd = 0x40}, а \texttt{payload}~--- это
\emph{упакованные} семплы сигнала.

\begin{figure}[htbp]
\centering
\adjustbox{max width=\textwidth}{%
\begin{tikzpicture}[
    field/.style={
        draw, thick,
        minimum height=1cm,
        align=center,
        font=\footnotesize\ttfamily,
        fill=white
    },
    meta/.style={field, fill=clrStateData!20},
    data/.style={field, fill=clrLogicData!20},
    label/.style={
        font=\footnotesize,
        above,
        inner sep=2pt
    },
    size/.style={
        font=\scriptsize,
        below,
        inner sep=2pt
    }
]

% Поля payload
\node[meta, minimum width=2cm] (fs) at (0, 0) {fs\_hz};
\node[meta, minimum width=1cm, right=0cm of fs] (ch) {ch};
\node[meta, minimum width=2cm, right=0cm of ch] (ns) {nsamples};
\node[meta, minimum width=1.5cm, right=0cm of ns] (pre) {pretrig};
\node[data, minimum width=4cm, right=0cm of pre] (samples) {samples[]};

% Подписи сверху
\node[label] at (fs.north) {u32};
\node[label] at (ch.north) {u8};
\node[label] at (ns.north) {u16};
\node[label] at (pre.north) {u16};
\node[label] at (samples.north) {u16 × nsamples};

% Размеры снизу
\node[size] at (fs.south) {4 байта};
\node[size] at (ch.south) {1};
\node[size] at (ns.south) {2};
\node[size] at (pre.south) {2};
\node[size] at (samples.south) {2 × N};

\end{tikzpicture}
}% end adjustbox
\caption{Структура \texttt{payload} кадра \texttt{OSC\_DATA}.}
\label{fig:protocol-osc-data}
\end{figure}

\subsection{Разбор полей}

\begin{itemize}
    \item \texttt{fs\_hz}~--- частота дискретизации, в герцах.
    \item \texttt{ch}~--- номер канала (у нас~--- \texttt{0}).
    \item \texttt{nsamples}~--- количество семплов в кадре.
    \item \texttt{pretrig}~--- сколько семплов \emph{до} триггера.
    \item \texttt{samples[]}~--- сами семплы (\texttt{u16}).
\end{itemize}

\begin{keybox}
\textbf{Почему \texttt{u16}, а не \texttt{u12}?} ADC у STM32F401~---
\SI{12}{\bit}. Но передавать \emph{упакованные} 12-битные
значения~--- \emph{сложно} (не выровнены по байтам).
Проще~--- \emph{дополнить} до 16 бит (\texttt{u16}).
Потеря~--- \emph{минимальная}, код~--- \emph{проще}.
\end{keybox}

\section{Отправка кадра}
\label{sec:protocol-send}

\textbf{Отправка}~--- это \emph{сборка} кадра и \emph{запись}
его в порт.

\begin{listing}[htbp]
\begin{minted}[frame=none, bgcolor=gray!5]{c}
static bool send_cmd(int fd, uint8_t cmd,
                     const uint8_t *payload,
                     uint16_t len, uint16_t *seq)
{
    if (fd < 0) return false;
    uint8_t buf[256];
    if (len > sizeof(buf) - 9) return false;

    uint16_t s = ++(*seq);
    size_t idx = 0;

    buf[idx++] = 0x55;       // sync low
    buf[idx++] = 0xAA;       // sync high
    buf[idx++] = 0x01;       // ver
    buf[idx++] = s & 0xFF;   // seq low
    buf[idx++] = s >> 8;     // seq high
    buf[idx++] = cmd;        // cmd
    buf[idx++] = len & 0xFF; // len low
    buf[idx++] = len >> 8;   // len high

    if (payload && len) {
        memcpy(&buf[idx], payload, len);
    }
    size_t hdr_len = idx + len;
    uint16_t crc = crc16_ibm(&buf[2], hdr_len - 2);
    buf[hdr_len] = crc & 0xFF;
    buf[hdr_len + 1] = crc >> 8;

    ssize_t w = write(fd, buf, hdr_len + 2);
    return w == (ssize_t)(hdr_len + 2);
}
\end{minted}
\caption{\texttt{send\_cmd()}~--- отправка кадра}
\label{lst:protocol-send}
\end{listing}

\subsection{Что здесь важно}

\begin{enumerate}
    \item \textbf{Буфер фиксированного размера} (\texttt{256}).
          \emph{Защита} от переполнения.
    \item \textbf{Инкремент} \texttt{seq}~--- \emph{каждый} кадр
          получает \emph{новый} номер.
    \item \textbf{Сборка} \emph{поля} за \emph{полем}~---
          \emph{байт} за \emph{байтом} (little-endian).
    \item \textbf{CRC} \emph{считается} от \texttt{ver} до конца
          \texttt{payload} (\texttt{buf[2]}~--- \emph{начало}).
    \item \textbf{Запись} \emph{всего} кадра за \emph{один}
          \texttt{write()}~--- \emph{атомарно}.
\end{enumerate}

\begin{warnbox}
\textbf{Почему \texttt{write()} за один раз?} Потому что это атомарно~--- целиком или никак. Если \emph{писать} по \emph{байту}, между вызовами может \emph{вклиниться} что-то ещё~--- \emph{нарушится} кадр. \emph{Один} \texttt{write()}~--- \emph{атомарная} операция.
\end{warnbox}

\begin{keybox}
\textbf{Что значит <<атомарно>>?} \emph{Атомарная} операция~---
это операция, которая выполняется \emph{целиком} или
\emph{не выполняется вовсе}. Никакой \emph{другой} поток
\emph{не может} вклиниться в её середину и увидеть
<<половину>>.

\textbf{Пример}: \texttt{write(fd, buf, 100)}~--- атомарно.
Ядро запишет \emph{все} 100 байт или \emph{ничего}.
А \texttt{100} вызовов \texttt{write(fd, buf+i, 1)}~---
\emph{не} атомарно: между ними \emph{может} вклиниться
другой поток и \emph{испортить} кадр.
\end{keybox}


\section{Приём кадра}
\label{sec:protocol-receive}

\textbf{Приём}~--- это \emph{самое} \emph{сложное}. Поток
байтов \emph{непрерывен}. \emph{Как} найти \emph{начало}
кадра?

\subsection{Пошаговый парсер}

\begin{listing}[htbp]
\begin{minted}[frame=none, bgcolor=gray!5]{c}
// Внутри osc_reader_thread
size_t pos = 0;
while (have >= 8) {
    // 1. Ищем sync 0x55 0xAA
    if (!(buf[pos] == 0x55 && buf[pos + 1] == 0xAA)) {
        pos++;
        have--;
        continue;
    }

    // 2. Сдвигаем буфер, если нужно
    if (pos > 0) {
        memmove(buf, buf + pos, have);
        pos = 0;
    }

    // 3. Читаем заголовок
    uint8_t ver = buf[2];
    uint16_t seq = buf[3] | (buf[4] << 8);
    uint8_t cmd = buf[5];
    uint16_t len = buf[6] | (buf[7] << 8);

    // 4. Ждём весь кадр
    size_t frame_len = 8 + len + 2;
    if (have < frame_len) break;

    // 5. Проверяем CRC
    uint16_t crc_calc = crc16_ibm(&buf[2], 6 + len);
    uint16_t crc_rx = buf[8 + len] | (buf[9 + len] << 8);
    if (crc_calc != crc_rx) {
        memmove(buf, buf + 2, have - 2);
        have -= 2;
        continue;
    }

    // 6. Обрабатываем кадр
    if (cmd == 0x40) {
        // OSC_DATA
        // ...
    }

    // 7. Сдвигаем буфер
    memmove(buf, buf + frame_len, have - frame_len);
    have -= frame_len;
}
\end{minted}
\caption{Парсер кадра в потоке}
\label{lst:protocol-receive}
\end{listing}

\subsection{Разбор шагов}

\begin{enumerate}
    \item \textbf{Поиск sync}. \emph{Сканируем} байты, пока
          не \emph{найдём} \texttt{0x55 0xAA}.
    \item \textbf{Сдвиг буфера}. \emph{Убираем} мусор \emph{до}
          sync.
    \item \textbf{Чтение заголовка}. \emph{Извлекаем} поля
          (little-endian).
    \item \textbf{Ожидание кадра}. Если в буфере \emph{меньше}
          \texttt{frame\_len}~--- \emph{ждём} следующего
          \texttt{read()}.
    \item \textbf{Проверка CRC}. Если \emph{не} совпала~---
          \emph{сдвигаем} на 2 байта (\emph{ищем} следующий sync).
    \item \textbf{Обработка}. \emph{Извлекаем} данные
          (для \texttt{OSC\_DATA}~--- семплы).
    \item \textbf{Сдвиг}. \emph{Убираем} обработанный кадр
          из буфера.
\end{enumerate}

\begin{keybox}
\textbf{Почему сдвиг на 2 байта при ошибке CRC?}
Потому что \emph{sync}~--- это 2 байта. Мы \emph{не знаем},
где \emph{следующий} sync. \emph{Сдвиг} на 2~---
\emph{быстрее}, чем на 1, и \emph{не} пропустит sync.
\end{keybox}

\begin{ideabox}
\textbf{Что такое <<конечный автомат>>?} \emph{Конечный
автомат} (Finite State Machine, FSM)~--- это модель
поведения, у которой есть \emph{конечный набор} состояний
и \emph{правила} переходов между ними.

\textbf{Пример}: светофор. Три состояния~--- красный,
жёлтый, зелёный. Переход~--- по таймеру. \emph{Нельзя}
быть <<полу-красным>>~--- либо одно, либо другое.

\textbf{Наш парсер~--- тоже автомат.} У него \emph{пять}
состояний:
\begin{enumerate}
    \item Поиск sync.
    \item Чтение заголовка.
    \item Ожидание полного кадра.
    \item Проверка CRC.
    \item Обработка кадра.
\end{enumerate}
Переходы~--- по событиям (пришёл байт, нашли sync,
CRC не совпала).

\textbf{Зачем это знать?} Если \emph{не} думать об
автомате~--- код \emph{зацикливается}, \emph{теряет}
состояние или \emph{падает}. Автомат \emph{дисциплинирует}:
ты \emph{чётко} знаешь, \emph{где} ты и \emph{что} делать.
\end{ideabox}

\begin{warnbox}
\textbf{Парсер~--- это \emph{конечный} автомат.} Если
написать \emph{неправильно}~--- буфер \emph{переполнится}
или \emph{зациклится}. \emph{Всегда} проверяйте \texttt{have}
и \emph{не} забывайте \texttt{memmove}.
\end{warnbox}

\subsection{Особенности \texttt{OSC\_DATA}}

Для кадра \texttt{OSC\_DATA} (\texttt{cmd = 0x40}) \emph{дополнительно}
проверяем:
\begin{itemize}
    \item \texttt{len >= 9} (минимум: \texttt{fs\_hz} + \texttt{ch} +
          \texttt{nsamples} + \texttt{pretrig}).
    \item \texttt{nsamples <= 2048} (защита от \emph{мусора}).
    \item \texttt{17 + nsamples * 2 <= 8 + len} (хватает ли
          данных).
\end{itemize}

\section{Уроки главы}
\label{sec:protocol-lessons}

\begin{enumerate}
    \item \textbf{Кадр~--- это контракт.} Формат
          \emph{фиксирован} и \emph{документирован}.
    \item \textbf{Sync~--- <<якорь>>.} Находит начало кадра
          в потоке.
    \item \textbf{Seq~--- защита от потерь и дубликатов.} Ответ
          содержит \emph{тот же} \texttt{seq}.
    \item \textbf{CRC16~--- контроль целостности.}
          \emph{Ловит} случайные ошибки. Если CRC не совпала~--- кадр \emph{выбрасываем}.
    \item \textbf{Частота в миллигерцах.} \emph{Нет}
          дробных чисел~--- \emph{проще} передавать.
    \item \textbf{Команды разделены по диапазонам:} генератор,
          осциллограф, данные.
    \item \textbf{OSC\_DATA~--- особый кадр.} \texttt{payload}~---
          семплы \texttt{u16}.
    \item \textbf{Отправка~--- \emph{один} write().}
          \emph{Атомарно}.
    \item \textbf{Приём~--- \emph{конечный} автомат.}
          \emph{Сканируем} sync, \emph{проверяем} CRC, \emph{сдвигаем}.
\end{enumerate}


\section{Чек-лист грабель}
\label{sec:protocol-pitfalls}

\begin{table}[htbp]
\centering
\caption{Грабли протокола}
\label{tab:protocol-pitfalls}
\small
\begin{tabularx}{\textwidth}{@{}c L L@{}}
\toprule
\textbf{№} & \textbf{Грабля} & \textbf{Решение} \\
\midrule
1 & Забыли sync & \texttt{0x55 0xAA} в начале \\
2 & Неверный CRC & Считать от \texttt{ver} до \texttt{payload} \\
3 & Буфер переполнен & Проверка \texttt{len > sizeof(buf) - 9} \\
4 & Потеря seq & Инкремент \emph{каждого} кадра \\
5 & Нет \texttt{memmove} после обработки & Сдвиг буфера \\
6 & Бесконечный цикл парсера & Проверка \texttt{have} \\
7 & Дробные числа & \emph{Миллигерцы} для частоты \\
8 & Мусор в начале потока & \emph{Сканируем} до sync \\
\bottomrule
\end{tabularx}
\end{table}

\section{Упражнения}
\label{sec:protocol-exercises}

\begin{exercisebox}
\textbf{Упражнение 1.} Вычислите вручную CRC16 для строки
\texttt{"Hello"} (5 байт). Проверьте через функцию
\texttt{crc16\_ibm()}.
\end{exercisebox}

\begin{exercisebox}
\textbf{Упражнение 2.} Нарисуйте \emph{диаграмму} кадра
\texttt{set\_freq} (\texttt{cmd=0x11}, \texttt{payload=4} байта).
Сколько байт всего? Где \texttt{crc}?
\end{exercisebox}

\needspace{5\baselineskip}
\begin{exercisebox}
\textbf{Упражнение 3.} Что произойдёт, если в кадре
\texttt{OSC\_DATA} \texttt{nsamples = 0}? Как парсер
должен \emph{отреагировать}?
\end{exercisebox}

\begin{exercisebox}
\textbf{Упражнение 4.} Добавьте в протокол \emph{новую}
команду \texttt{set\_filter} (\texttt{cmd=0x16}) с
\texttt{payload} из \emph{двух} байт (коэффициенты).
\emph{Опишите} формат.
\end{exercisebox}

\section{Проверьте сами}
\label{sec:protocol-check}

\begin{enumerate}
    \item Откройте \texttt{docs/protocol.md} в репозитории.
    \item Сравните с тем, что написано в главе.
    \item Найдите \emph{несоответствия} (если есть).
    \item Соберите \emph{вручную} кадр \texttt{set\_wave(0)}~---
          синус (\texttt{cmd=0x10}, \texttt{payload=[0]}).
    \item Вычислите CRC \emph{вручную}.
    \item Сравните с тем, что \emph{отправит} \texttt{send\_cmd()}.
\end{enumerate}

\begin{ideabox}
\textbf{Если получилось~--- вы поняли протокол.} Это
\emph{главное}. Протокол~--- \emph{фундамент} общения.
Без него~--- \emph{нет} осциллографа, \emph{нет} генератора.
\end{ideabox}

